% ==================== 附录 ====================
\appendix
\ctexset{chapter/format = \centering\zihao{3}\heiti\bfseries,
         chapter/name = {附录,~},
         chapter/number = \Alph{chapter}}

\chapter{Apollo Planning 全流程数值仿真程序源码}
\label{appendix:apollo-pipeline}

\zihao{-4}\songti
本附录列出第\ref{chap:integration}章数值仿真所用程序的完整源码。该程序按 Apollo \texttt{lane\_follow} 流水线顺序复现路径与速度规划的七个核心阶段，Baseline 与 MIKU 两种算法在同一文件内并列实现，所需依赖通过 PEP~723 内联元数据声明。

\newfontfamily\codefont{Hack}[Scale=0.95]
\lstinputlisting[
  language=Python,
  basicstyle=\linespread{0.95}\selectfont\zihao{6}\codefont,
  keywordstyle=\color{deepblue},
  commentstyle=\color{lightgreen!60!black},
  stringstyle=\color{dangerred!60!black},
  numbers=left,
  numberstyle=\tiny\color{gray},
  frame=single,
  breaklines=true,
  breakatwhitespace=false,
  columns=flexible,
  showstringspaces=false,
  tabsize=4,
  extendedchars=true,
  inputencoding=utf8,
  aboveskip=4pt,
  belowskip=4pt,
]{../可视化/apollo_pipeline.py}

\chapter{MIKU 路径边界决策器 C++ 实现}
\label{appendix:path-bounds}

\zihao{-4}\songti
本附录列出 MIKU 算法在 Apollo 框架中的核心 C++ 实现。\texttt{PathBoundsDeciderUtil} 类承载了扫描线分组、交互密度威胁度计算、max-gap 最优通行间隙选择等全部路径边界决策逻辑。

\vspace{6pt}
\noindent\textbf{path\_bounds\_decider\_util.h}

\lstinputlisting[
  language=C++,
  basicstyle=\linespread{0.95}\selectfont\zihao{6}\codefont,
  keywordstyle=\color{deepblue},
  commentstyle=\color{lightgreen!60!black},
  stringstyle=\color{dangerred!60!black},
  numbers=left,
  numberstyle=\tiny\color{gray},
  frame=single,
  breaklines=true,
  breakatwhitespace=false,
  columns=flexible,
  showstringspaces=false,
  tabsize=4,
  aboveskip=4pt,
  belowskip=4pt,
]{code/path_bounds_decider_util.h}

\vspace{6pt}
\noindent\textbf{path\_bounds\_decider\_util.cc}

\lstinputlisting[
  language=C++,
  basicstyle=\linespread{0.95}\selectfont\zihao{6}\codefont,
  keywordstyle=\color{deepblue},
  commentstyle=\color{lightgreen!60!black},
  stringstyle=\color{dangerred!60!black},
  numbers=left,
  numberstyle=\tiny\color{gray},
  frame=single,
  breaklines=true,
  breakatwhitespace=false,
  columns=flexible,
  showstringspaces=false,
  tabsize=4,
  aboveskip=4pt,
  belowskip=4pt,
]{code/path_bounds_decider_util.cc}

\chapter{MIKU 车道跟随路径任务与速度优化器实现}
\label{appendix:lane-speed}

\zihao{-4}\songti
本附录列出 MIKU 改动涉及的车道跟随路径任务（集成 MIKU 边界决策）与分段加加速度速度优化器（ST 走廊裁剪与 QP 求解）的完整实现。

\vspace{6pt}
\noindent\textbf{lane\_follow\_path.cc}

\lstinputlisting[
  language=C++,
  basicstyle=\linespread{0.95}\selectfont\zihao{6}\codefont,
  keywordstyle=\color{deepblue},
  commentstyle=\color{lightgreen!60!black},
  stringstyle=\color{dangerred!60!black},
  numbers=left,
  numberstyle=\tiny\color{gray},
  frame=single,
  breaklines=true,
  breakatwhitespace=false,
  columns=flexible,
  showstringspaces=false,
  tabsize=4,
  aboveskip=4pt,
  belowskip=4pt,
]{code/lane_follow_path.cc}

\vspace{6pt}
\noindent\textbf{piecewise\_jerk\_speed\_optimizer.cc}

\lstinputlisting[
  language=C++,
  basicstyle=\linespread{0.95}\selectfont\zihao{6}\codefont,
  keywordstyle=\color{deepblue},
  commentstyle=\color{lightgreen!60!black},
  stringstyle=\color{dangerred!60!black},
  numbers=left,
  numberstyle=\tiny\color{gray},
  frame=single,
  breaklines=true,
  breakatwhitespace=false,
  columns=flexible,
  showstringspaces=false,
  tabsize=4,
  aboveskip=4pt,
  belowskip=4pt,
]{code/piecewise_jerk_speed_optimizer.cc}
